\subsubsection{Représentation pentadique de la microfibre et de son incidence}
\label{pi-estrella:desarrollo}

La disposition des cinq régions permet de visualiser simultanément
l'unité du quotient ternaire et la différenciation conservée par le
catalogue. L'application pertinente est la restriction
\[
 \Psi\big|_{\mathscr R_\pi^{\rm micro}}:
 \mathscr R_\pi^{\rm micro}\longrightarrow\{010211\},
 \qquad \Psi(U)=(-U)\bmod3.
\]
Ses cinq antécédents sont les vecteurs \(U_6\) du tableau précédent. La figure \ref{pi-estrella:figura} restitue leur organisation radiale, en conservant l'ordre de coordinatisation, les signatures et les multiplicités du catalogue (équation~\eqref{eq:palabras-regionales}).

% Orden de carta y datos recuperados de fig_cinco_regiones_pi.tex.
\begin{figure}[H]
\centering
\begin{tikzpicture}[x=1cm,y=1cm,font=\small,>=Latex]
 \coordinate (a) at (90:2.2);
 \coordinate (b) at (18:2.2);
 \coordinate (c) at (-54:2.2);
 \coordinate (d) at (-126:2.2);
 \coordinate (e) at (162:2.2);
 \draw[black!35,line width=.45pt]
    (a)--(c)--(e)--(b)--(d)--cycle;
 \node[draw,fill=white,circle,minimum size=1.6cm,align=center,
    inner sep=3pt] (z) at (0,0) {$w_\pi$\\$010211$};
 \foreach \v in {a,b,c,d,e}{
    \fill (\v) circle (2pt);
    \draw[->,line width=.6pt] (\v)--(z);
 }
 \node[above=3mm,align=center] at (a)
   {$U_\pi^\perp$\\$555555$\\multiplicidad $432$};
 \node[right=3mm,align=left] at (b)
   {$U_{\pi,1}^{++}$\\$339555$\\multiplicidad $144$};
 \node[below right=2mm and 1mm,align=left] at (c)
   {$U_{\pi,2}^{++}$\\$393555$\\multiplicidad $144$};
 \node[below left=2mm and 1mm,align=right] at (d)
   {$U_{\pi,3}^{++}$\\$933555$\\multiplicidad $144$};
 \node[left=3mm,align=right] at (e)
   {$U_\pi^{--}$\\$933717$\\multiplicidad $144$};
 \node[align=center,font=\footnotesize] at (0,-3.55)
   {$1_\perp+3_{++}+1_{--}$\qquad $432+4\cdot144=1008$};
\end{tikzpicture}
\caption{Les cinq régions de \(\pi\) en disposition pentadique.
Chaque position conserve sa signature multiplicative et sa multiplicité ;
les flèches radiales représentent la projection commune
\(\Psi(U)=010211\). Le contour étoilé organise les cinq
positions de la coordinatisation. Leur réalisation dans les cinq octades de Witt
s'effectue par l'alignement régional et la décomposition
\(4+1\) du texte.}
\label{pi-estrella:figura}
\end{figure}


La projection sur le centre identifie les cinq images ternaires ; le
registre des positions extérieures retient l'information régionale
qui rend possible la réalisation d'incidence. La distribution
\(1_\perp+3_{++}+1_{--}\) spécifie une région transversale, trois
régions positives et une négative. Leurs poids normalisés sont
\(3/7\) et \(1/7\) pour chacune des quatre autres : la
multiplicité transversale est triple, alors même que les cinq images par
\(\Psi\) coïncident.

\Needspace{11\baselineskip}
Le lien avec l'étoile de Witt se formule au moyen du drapeau
\(B_K\subset H^-_\alpha\), de l'alignement \(\ell\) et du
relèvement \(\iota_D\) déjà construits. Les cinq positions régionales
se réalisent dans les cinq octades sur \(\ell(B_K)\) :
\[
\begin{aligned}
 U_\pi^\perp&\longmapsto O^\perp_{\ell(B_K)},\\
 U_\pi^{--}&\longmapsto\iota_D(H^-_\alpha),\\
 \{U_{\pi,1}^{++},U_{\pi,2}^{++},U_{\pi,3}^{++}\}
 &\longmapsto
 \{\iota_D(B_K\cup P):P\in\mathcal P_+\},\\
 \mathcal P_+&=\bigl\{\{1,3\},\{2,11\},\{8,12\}\bigr\}.
\end{aligned}
\]
L'assignation des trois positions positives aux trois paires
s'effectue dans une coordinatisation ordonnée. La ligne transversale et la ligne négative
sont distinguées par leurs prédicats régionaux et par le drapeau.

Les deux figures s'articulent ainsi : la figure
\ref{exc:fig-estrella} montre les quatre hexades du petit plan
de Witt sur \(B_K\) ; la figure \ref{pi-estrella:figura} montre
les cinq régions qui, sous l'alignement, occupent les quatre octades
résiduelles et l'octade transversale du plan binaire. La loi
\(4+1\), démontrée dans la proposition
\ref{exc:residuo-binario}, spécifie le changement d'incidence. Dans cette
composition, \(K\) détermine la tétrade marquée et la signature de
\(\alpha\) distingue l'une de ses hexades. Les régions conservent,
outre leur valeur ternaire commune, la structure nécessaire à ce
transport vers la réalisation exceptionnelle.
